\documentclass[12pt,a4paper]{article}

\usepackage[margin=25mm]{geometry}
\usepackage{newtxtext,newtxmath}
\usepackage{microtype}
\usepackage{setspace}
\usepackage{siunitx}
\usepackage{booktabs}
\usepackage{tabularx}
\usepackage{enumitem}
\usepackage{xurl}
\usepackage[hidelinks]{hyperref}
\usepackage[nameinlink,noabbrev]{cleveref}

\onehalfspacing
\setlength{\parindent}{1.25em}
\setlength{\parskip}{0.25em}
\sisetup{detect-all}

\title{Analysis, Design, and Fabrication of a Multimodal Instrumented Prosthetic-Liner Slab}
\author{}
\date{}

\begin{document}
\maketitle

\section{Introduction and Background}
\label{sec:introduction}

% To be developed after the literature review has fixed the terminology,
% research gap, and defensible scope of the project.

\section{Literature Review}
\label{sec:literature-review}

\subsection{The prosthetic socket--residual-limb interface as a sensing problem}

The mechanical and thermal conditions at the interface between a residual limb,
liner, and prosthetic socket strongly influence comfort and tissue health.
The liner must simultaneously suspend the prosthesis, distribute load, limit
relative motion, and protect soft tissue. These functions create a difficult
measurement environment. Normal pressure is produced as load is transferred
between the limb and socket, while tangential motion produces shear. The liner
also forms a comparatively closed thermal and moisture barrier. Heat and
perspiration can therefore accumulate during activity, changing both the
condition of the skin and the mechanical behaviour of the interface. A useful
instrumented liner should consequently measure more than one isolated variable.
Normal pressure, shear, temperature, and humidity represent complementary
aspects of the same physical problem.

The importance of this combined view is supported by the broader interface-
measurement literature. Al-Fakih \textit{et al.} reviewed techniques used to
measure socket-interface stresses and showed that pressure has historically
received more attention than shear, even though both contribute to the
mechanical state of the residuum \cite{alfakih2016}. Commercial pressure films
and discrete transducers can provide useful measurements, but their thickness,
wiring, calibration, and sensitivity to curvature can disturb the interface
being measured. Shear measurement is more demanding because a sensor must
separate tangential loading from normal compression while remaining sufficiently
thin and compliant. This is a central design constraint for the present work:
the sensing system must not create a local prominence that changes the load
distribution or becomes a new source of discomfort.

Thermal and moisture measurements introduce a related set of constraints.
Klute \textit{et al.} demonstrated that prosthetic socket and liner materials
form barriers to heat transfer \cite{klute2007}. The consequences are not only
subjective warmth. Increased temperature promotes perspiration, while trapped
moisture may alter friction, reduce suspension, soften the skin, and contribute
to irritation. Paternò \textit{et al.} accordingly argue that temperature and
humidity should be studied alongside mechanical loading rather than treated as
secondary comfort variables \cite{paterno2020}. The literature therefore
supports a multimodal architecture; however, it also shows that the utility of
such an architecture depends on integration. A sensor that is accurate in
isolation is not necessarily suitable when embedded in a flexible liner.

\subsection{Multiaxial capacitive sensing: Albrecht \textit{et al.}}

Albrecht \textit{et al.} developed an inkjet-printed capacitive sensor intended
to distinguish normal and two-directional shear loading
\cite{albrecht2019}. The design is especially relevant because it addresses
the cross-sensitivity problem without relying on a thick, rigid multi-axis load
cell. Four capacitors, denoted \(C_1\)--\(C_4\), are formed by patterned
conductive layers separated by a compliant dielectric. The electrodes use an
E-shaped geometry. Multiple electrode ``fingers'' increase the edge length over
which overlap changes when one conductive layer moves relative to the other.
The initial overlap is arranged near the middle of its working range, allowing
the capacitance of opposing elements to increase or decrease with the direction
of shear.

The four-capacitor arrangement provides a physical basis for resolving the
loading components. Under compression, the dielectric thickness changes and
the capacitors exhibit a broadly common-mode response. Under shear in one
direction, one opposed capacitor pair responds differentially; the orthogonal
pair performs the same function for the second tangential direction.
Consequently, sums and differences of the four outputs can be used to estimate
normal force and signed shear in two axes. This is a more appropriate model of
the prosthetic interface than a pressure-only sensor because socket loading is
not purely normal. Pistoning, rotation, and gait-related relative motion can
produce tangential forces even where normal pressure appears acceptable.

The reported device used inkjet-deposited conductive traces approximately
\SI{412}{\nano\metre} thick and a compliant dielectric approximately
\SI{40}{\micro\metre} thick. The electrode fingers were approximately
\SI{210}{\micro\metre} wide with a \SI{40}{\micro\metre} gap, and each
finger was approximately \SI{8}{\milli\metre} long. Over an investigated range
of approximately \SIrange{0.1}{8}{\newton}, Albrecht \textit{et al.} reported a
normal-force sensitivity of about \SI{5.2}{\femto\farad\per\newton} and a
shear-force sensitivity of about
\SI{13.1}{\femto\farad\per\newton} \cite{albrecht2019}. The differential shear
response was direction-dependent, while the four-element arrangement could
increase the combined output by suitable signal aggregation.

These findings make the architecture a credible starting point for the slab
sensor, but they do not by themselves establish performance inside a prosthetic
liner. Several translation questions remain. First, the force range and
boundary conditions of the published test must be compared with the pressures,
contact areas, and shear forces expected in the planned slab test. Second,
capacitance at the femtofarad-per-newton scale is vulnerable to cable,
connector, and environmental parasitics. Third, encapsulation changes the
effective stiffness and therefore the mechanical transfer of force to the
dielectric. Finally, repeated flexing, sweat, and temperature variation can
change the baseline or produce drift. The value of Albrecht's study is thus
twofold: it supplies a defensible transducer geometry and it identifies the
calibration problem that the present project must solve after embedding.

\subsection{Printed capacitive humidity sensing: Aït-Mammar \textit{et al.}}

Aït-Mammar \textit{et al.} investigated an inkjet-fabricated capacitive humidity
sensor formed on a flexible polyimide substrate
\cite{aitmammar2020}. Their approach used interdigitated electrodes covered by
a hygroscopic dielectric, cellulose acetate butyrate (CAB). Interdigitated
electrodes are attractive for a liner because both terminals can be patterned
on one surface; a separate parallel counter-electrode is unnecessary. The
electric field extends through the CAB layer and the surrounding environment.
As the polymer absorbs water, its effective permittivity changes, producing a
measurable change in capacitance.

The cited device employed 52 electrode pairs with approximately
\SI{100}{\micro\metre} gaps and a CAB sensing layer approximately
\SI{8}{\micro\metre} thick. The conductive layer was on the order of
\SI{300}{\nano\metre}, while the Kapton substrate was approximately
\SI{125}{\micro\metre} thick. The study characterised the sensor over a broad
relative-humidity range, approximately \(20\)--\(90\%\), and reported a clear
humidity-dependent capacitance response \cite{aitmammar2020}. This construction
is relevant to the proposed slab because it combines a thin substrate, printed
conductors, and a solution-processable sensing film. In principle, four small
elements can be distributed across the slab to measure spatial rather than
single-point moisture conditions.

The literature also exposes a fundamental integration conflict. A humidity
sensor must exchange water vapour with the measured environment. Complete
encapsulation in impermeable silicone would protect the conductors but would
also isolate the CAB layer and suppress or delay its response. The sensing
surface therefore needs a vapour-access path, selective membrane, or local
exposure to the liner--skin environment. That access route must prevent liquid
sweat from shorting the electrodes or causing uncontrolled swelling. Response
and recovery times measured in a controlled humidity chamber cannot be assumed
to apply under occlusion, contact pressure, condensation, or repeated wetting.
Likewise, the published operating frequency and capacitance range must inform
the readout circuit; they cannot simply be transferred to an arbitrary
capacitance-to-digital converter.

For the present work, Aït-Mammar's study justifies the material principle and
planar geometry, not a finished prosthetic sensor. The required follow-on work
is to verify CAB deposition on the available printed electrode, characterise
film thickness and uniformity, determine whether the proposed opening or
membrane preserves vapour access, and calibrate the completed element at
skin-relevant temperatures. Hysteresis, saturation, recovery after exposure to
high humidity, and cross-sensitivity to bending should be reported. These tests
are necessary before the humidity elements can be interpreted as measurements
of local liner microclimate.

\subsection{Printed resistive temperature sensing: Liew \textit{et al.}}

Liew and Lee studied inkjet-printed resistive temperature detectors fabricated
from silver nanoparticle inks on flexible polyethylene terephthalate (PET)
\cite{liew2020}. The transducer uses the temperature coefficient of electrical
resistance: as temperature rises, the resistance of the printed conductor
changes. A serpentine geometry increases conductor length within a small area,
raising the nominal resistance and measurable temperature-dependent change
without requiring a large sensor footprint. The reported pattern was
approximately \SI{6.9}{\milli\metre} by \SI{13.2}{\milli\metre}, with a line
width near \SI{0.5}{\milli\metre}, spacing near
\SI{0.4}{\milli\metre}, and \SI{1.5}{\milli\metre} contact pads.

The study compared two silver nanoparticle inks over approximately
\SIrange{30}{100}{\degreeCelsius}. A linear relationship between resistance and
temperature was reported for both printed conductors. The reported sensitivity
was approximately \SI{0.1086}{\ohm\per\degreeCelsius} for one ink and
\SI{0.0543}{\ohm\per\degreeCelsius} for the JS-B25HV ink
\cite{liew2020}. The latter is particularly relevant where the same commercial
ink and a Fujifilm Dimatix printer are available. The study therefore provides
a direct process precedent for printing four temperature elements on a thin
polymer substrate.

Nevertheless, a linear result over \SIrange{30}{100}{\degreeCelsius} is not a
complete calibration for prosthetic use. The most important physiological
changes may occur over a much narrower interval near skin temperature. The
readout must resolve small resistance changes within that interval while
separating them from lead resistance, contact resistance, mechanical strain,
and self-heating. Silver traces on PET are not mechanically neutral: bending
and tensile strain can change resistance or initiate cracking, making a nominal
temperature signal partly strain-dependent. Encapsulation can improve
environmental durability but adds thermal resistance and time lag.

Liew and Lee consequently support the feasibility of the printed RTD, while the
present slab must establish application-specific metrology. Each element should
be calibrated after printing and again after embedding, against a traceable
reference over a skin-relevant range. Thermal response time, repeatability,
drift, bending sensitivity, and agreement between the four elements should be
reported. A four-wire measurement would reduce lead-resistance error where
space and interconnect count permit; otherwise, conductor and lead resistance
must be included in the calibration model.

\subsection{Discrete temperature measurement and thermal control:
Ghoseiri \textit{et al.}}

Ghoseiri \textit{et al.} developed a temperature measurement and control system
for transtibial prostheses using 12 temperature sensors distributed on opposing
sides of the interface \cite{ghoseiri2018}. Their work connects temperature
measurement to a clinically meaningful mechanism: heat accumulation promotes
perspiration, and the combined warm and humid environment can affect comfort,
skin condition, and prosthetic suspension. The study also shows the value of
spatial measurements. A single temperature value cannot represent gradients
created by local loading, anatomy, material thickness, and nonuniform heat
transfer.

The system did more than log temperature. Measured differences were used to
control heating or cooling, demonstrating a closed-loop thermal-management
concept. Its complexity, however, also reveals the practical limits of adding
components to a wearable socket. Sensors, conductive paths, control
electronics, and thermal-management hardware add stiffness, volume, power
demand, and possible failure points. Ghoseiri \textit{et al.} identified further
miniaturisation and system refinement as prerequisites for practical
translation \cite{ghoseiri2018}.

The present slab uses two TMP36 integrated-circuit sensors as discrete internal
references. According to the manufacturer, the TMP36 operates from
\SIrange{2.7}{5.5}{\volt}, has a nominal output slope of
\SI{10}{\milli\volt\per\degreeCelsius}, and is specified over
\SIrange{-40}{125}{\degreeCelsius} \cite{tmp36}. Its simple voltage output and
low current make it convenient for prototype acquisition. Two sensors can
indicate whether a thermal gradient exists across the slab and can provide an
independent comparison with the printed RTDs.

The TMP36 choice is experimental rather than automatically liner-compatible.
The TO-92 package is substantially thicker and stiffer than a printed trace.
Within a \SI{3.5}{\milli\metre} middle layer, a package approximately
\SI{2.8}{\milli\metre} thick leaves little material above and below it. Local
stiffness and stress concentration must therefore be assessed. The sensors
measure their package temperature, not instantaneous skin temperature, and the
surrounding silicone introduces lag. Ghoseiri's work supports distributed
temperature measurement, while the TMP36 datasheet supports the electrical
choice; neither source proves that this package will remain comfortable or
mechanically safe inside a final liner.

\subsection{An instrumented liner precedent: Paternò \textit{et al.}}

Paternò \textit{et al.} provide the closest system-level precedent for the
present project \cite{paterno2020}. They developed a custom transtibial liner
with 12 embedded Hygrochron temperature-and-humidity loggers and compared its
microclimate measurements with those obtained using a participant's
conventional liner. The liner was designed from the user's geometry, and
sensor locations were incorporated through a controlled manufacturing route
that included machining. The work demonstrates that distributed temperature
and relative-humidity measurement can be integrated into a wearable liner and
used during a structured protocol rather than only in a benchtop chamber.

The study is important for three reasons. First, it frames temperature and
humidity as spatially varying quantities. Paternò \textit{et al.} reported mean
conditions at the start of one protocol of approximately
\SI{31.42}{\degreeCelsius} and \(68.13\%\) relative humidity, with values changing
during rest and activity \cite{paterno2020}. Second, it demonstrates a route
from user geometry to instrumented liner manufacture. Third, it acknowledges
that the sensor must be integrated without undermining the liner's primary
mechanical and comfort functions.

Its limitations are equally instructive. The Hygrochron devices are rigid
packages and are much larger than printed thin-film elements. The work was an
early case study involving a limited participant sample and a controlled
protocol. It does not establish long-term durability, washability, or universal
comfort. Nor does it measure normal pressure or shear, so the recorded
microclimate cannot be directly related to local mechanical loading. These
limitations define an opportunity for the present design: printed elements may
reduce thickness and increase spatial resolution, while a multiaxial
capacitive element can add the mechanical variables absent from the
instrumented-liner precedent.

The proposed slab should not be presented as a direct miniaturisation of
Paternò's liner. It combines sensing mechanisms with different substrates,
readout requirements, environmental sensitivities, and mechanical stiffnesses.
The slab is instead an integration experiment intended to determine whether
those mechanisms can coexist within a silicone laminate. A flat slab is an
appropriate preliminary specimen because layer placement, encapsulation,
electrical continuity, cross-sensitivity, and local compliance can be examined
before undertaking the more expensive and less repeatable manufacture of a
full curved liner.

\subsection{Synthesis and research gap}

No single reviewed study supplies the complete solution required here.
Albrecht \textit{et al.} provide a thin capacitive architecture capable of
separating normal and tangential loading, but not a validated prosthetic-liner
implementation. Aït-Mammar \textit{et al.} provide a printable humidity-sensing
mechanism, but integration must preserve vapour access. Liew and Lee provide a
printed silver RTD process, but it requires calibration near skin temperature
and separation of thermal and strain effects. Ghoseiri \textit{et al.} show why
distributed thermal measurement matters, but also illustrate the bulk and
complexity of active socket systems. Paternò \textit{et al.} demonstrate that
temperature and humidity can be logged from an instrumented liner during use,
but their rigid packages do not provide pressure or shear data and limit
miniaturisation.

The resulting research gap is not merely the absence of another sensor. It is
the absence of a compact, mechanically integrated test article that combines:
\begin{enumerate}[label=(\roman*),leftmargin=2em]
    \item directional mechanical sensing;
    \item distributed printed temperature measurement;
    \item distributed printed humidity measurement;
    \item discrete temperature references; and
    \item a silicone layer structure representative of a liner.
\end{enumerate}
The proposed \SI{56}{\milli\metre} by \SI{55}{\milli\metre} by
\SI{5.2}{\milli\metre} slab addresses this gap at prototype scale. Its three
layers comprise a \SI{0.20}{\milli\metre} first layer, a
\SI{3.50}{\milli\metre} sensor-containing middle layer, and a
\SI{1.50}{\milli\metre} cover layer. The planned sensing set comprises one
four-capacitor normal/shear element, two TMP36 devices, four printed RTDs, and
four printed humidity elements.

This literature supports the architecture as a research hypothesis, not as a
validated result. The decisive contribution must come from fabrication and
testing. The completed slab must be evaluated for dimensional fidelity,
electrical continuity, local stiffness, normal/shear decoupling, temperature
accuracy, humidity response, response time, repeatability, drift, hysteresis,
and behaviour after bending or cyclic loading. Only then can the design be
judged suitable for progression from a flat integration coupon to a curved
prosthetic liner.

\section{Methodology}
\label{sec:methodology}

% To be written only after the actual materials, fabrication sequence,
% instrumentation, calibration protocol, and analysis plan have been confirmed.

\section{Results}
\label{sec:results}

% No results should be inserted until measurements have been obtained.

\section{Discussion and Conclusion}
\label{sec:discussion}

% To be developed after the results are available.

\begin{thebibliography}{99}

\bibitem{alfakih2016}
E.~A. Al-Fakih, N.~A. Abu Osman, and F.~R. Mahmad Adikan,
``Techniques for interface stress measurements within prosthetic sockets of
transtibial amputees: A review of the past 50 years,''
\textit{Sensors}, vol.~16, no.~7, art.~1119, 2016.
\href{https://doi.org/10.3390/s16071119}{doi:10.3390/s16071119}.

\bibitem{klute2007}
G.~K. Klute, G.~I. Rowe, A.~V. Mamishev, and W.~R. Ledoux,
``The thermal conductivity of prosthetic sockets and liners,''
\textit{Prosthetics and Orthotics International}, vol.~31, no.~3,
pp.~292--299, 2007.
\href{https://doi.org/10.1080/03093640601042554}
{doi:10.1080/03093640601042554}.

\bibitem{albrecht2019}
A.~Albrecht, M.~Trautmann, M.~Becherer, P.~Lugli, and A.~Rivadeneyra,
``Shear-force sensors on flexible substrates using inkjet printing,''
\textit{Journal of Sensors}, vol.~2019, art.~1864239, 2019.
\href{https://doi.org/10.1155/2019/1864239}
{doi:10.1155/2019/1864239}.

\bibitem{aitmammar2020}
W.~Aït-Mammar, S.~Bridonneau, V.~Noël, E.~Stavrinidou, B.~Piro, and
G.~Mattana,
``A flexible capacitive humidity sensor based on inkjet-printed electrodes
and a cellulose acetate butyrate sensing layer,''
2020.
\href{https://doi.org/10.3390/ecsa-7-08216}
{Source details to be verified against the original project PDF before final
submission.}

\bibitem{liew2020}
Q.~J. Liew and H.~W. Lee,
``Inkjet-printed resistive temperature sensor based on silver nanoparticle
ink,''
in \textit{Proceedings of the 7th International Electronic Conference on
Sensors and Applications}, art.~8216, 2020.
\href{https://doi.org/10.3390/ecsa-7-08216}
{doi:10.3390/ecsa-7-08216}.

\bibitem{ghoseiri2018}
K.~Ghoseiri, Y.~Zheng, K.~L. Leung, M.~Rahgozar, G.~Aminian,
T.~H. Lee, and M.~R. Safari,
``Temperature measurement and control system for transtibial prostheses,''
\textit{Assistive Technology}, vol.~30, no.~1, pp.~16--23, 2018.
\href{https://doi.org/10.1080/10400435.2016.1225850}
{doi:10.1080/10400435.2016.1225850}.

\bibitem{paterno2020}
L.~Paternò, V.~Dhokia, A.~Menciassi, J.~Bilzon, and E.~Seminati,
``A personalised prosthetic liner with embedded sensor technology,''
\textit{BioMedical Engineering OnLine}, vol.~19, art.~71, 2020.
\href{https://doi.org/10.1186/s12938-020-00814-y}
{doi:10.1186/s12938-020-00814-y}.

\bibitem{tmp36}
Analog Devices,
``TMP35/TMP36/TMP37 low voltage temperature sensors,'' Rev.~H,
data sheet, 2021.
\url{https://www.analog.com/media/en/technical-documentation/data-sheets/TMP35_36_37.pdf}.

\end{thebibliography}

\end{document}
